\documentclass[11pt]{article}
\usepackage[margin=1in]{geometry}
\usepackage{amsmath,amssymb,amsthm,bm}
\usepackage{xcolor}
\usepackage{hyperref}
\usepackage{microtype}
\hypersetup{colorlinks=true,linkcolor=blue!50!black,citecolor=blue!50!black,urlcolor=blue!50!black}

\newtheorem{proposition}{Proposition}[section]
\newtheorem{lemma}[proposition]{Lemma}
\newtheorem{corollary}[proposition]{Corollary}
\newtheorem{theorem}[proposition]{Theorem}
\theoremstyle{definition}
\newtheorem{definition}[proposition]{Definition}
\newtheorem{remark}[proposition]{Remark}
\newtheorem{exercise}[proposition]{Exercise}

\newcommand{\q}{\mathbf{q}}
\newcommand{\x}{\mathbf{x}}
\newcommand{\kk}{\mathbf{k}}
\newcommand{\Psic}{\Psi_{c}}
\newcommand{\Psif}{\Psi_{f}}
\newcommand{\PW}{\mathcal{P}_W}
\newcommand{\Bf}{\mathcal{B}_f}
\newcommand{\Bc}{\mathcal{B}_c}
\newcommand{\BW}{\mathcal{B}_W}
\newcommand{\E}{\mathbb{E}}
\newcommand{\R}{\mathbb{R}}
\newcommand{\Z}{\mathbb{Z}}
\newcommand{\dd}{\mathrm{d}}
\newcommand{\tr}{\operatorname{tr}}
\newcommand{\Var}{\operatorname{Var}}
\newcommand{\Ny}{{\rm Ny}}
\newcommand{\lin}{{\rm lin}}
\newcommand{\Zel}{\mathcal{Z}}
\newcommand{\Nbody}{\mathcal{S}}

\title{Progressive $2\times$ super-resolution of Lagrangian $N$-body fields:\\ a self-contained derivation of the ``octave flow'' method}
\author{Notes prepared for Xiaowen Zhang}
\date{September 8, 2026 (v2, revised after review)}

\begin{document}
\maketitle

\begin{abstract}
\noindent These notes derive, from first principles, every mathematical statement behind the method implemented in \texttt{phase0\_octaves.py} and \texttt{octave\_flow\_toy.py}: the Lagrangian multi-resolution setup with nested initial conditions; the restriction/prolongation pair $(R,P)$ and the resulting orthogonal decomposition of a fine field into a coarse part and a detail part; the linear-theory statements (the coarse correction vanishes at first order, the detail is an independent Gaussian, the dimension of the new randomness); the second-order Lagrangian perturbation theory (2LPT) split that gives the $k^{2}$/$k^{4}$ suppression of the coarse correction and shows that the octave couples to the coarse field through its deformation tensor; the exact push-forward and Markov properties of the conditional distribution; the conditional flow-matching theory, including why the physical coupling is admissible and why one-step regression loses power; the scale-free self-similarity used for weight sharing; and the evaluation metrics. Exercises mark the places where you may want to redo the algebra yourself.
\end{abstract}

\tableofcontents

%======================================================================
\section{Setting and conventions}
%======================================================================

\subsection{Lagrangian fields on nested grids}

The simulation box is the torus $\mathbb{T}^3=[0,L)^3$ with periodic boundary conditions. A resolution level $\ell$ has $N_\ell$ particles per dimension, grid spacing $h_\ell=L/N_\ell$, and Lagrangian grid
\begin{equation}
\Lambda_\ell=\{\q_{\mathbf i}=(\mathbf i+o\,\mathbf 1)\,h_\ell:\ \mathbf i\in\{0,\dots,N_\ell-1\}^3\},\qquad o\in\{0,\tfrac12\},
\end{equation}
where $o$ is the grid offset used by the initial-condition generator. Two successive levels satisfy $N_{\ell+1}=2N_\ell$, $h_{\ell+1}=h_\ell/2$. In what follows a generic pair of levels is called coarse ($c$) and fine ($f$).

A particle that starts at $\q$ is at $\x(\q,t)=\q+\Psi(\q,t)$ at time $t$. The displacement $\Psi:\Lambda_\ell\to\R^3$ (and the velocity $\mathbf v$) are the objects we super-resolve. They are stored as three-channel cubes; all of the analysis treats them as periodic vector fields sampled on $\Lambda_\ell$.

\subsection{Fourier conventions}

For a periodic function $g$ on the torus,
\begin{equation}
\hat g(\kk)=\int_{\mathbb T^3} g(\q)\,e^{-i\kk\cdot\q}\,\dd^3q,\qquad g(\q)=\frac{1}{L^3}\sum_{\kk\in k_F\Z^3}\hat g(\kk)\,e^{i\kk\cdot\q},\qquad k_F=\frac{2\pi}{L}.
\label{eq:fourier}
\end{equation}
Products become convolutions with the factor $L^{-3}$:
\begin{equation}
\widehat{g_1g_2}(\kk)=\frac{1}{L^3}\sum_{\kk_1+\kk_2=\kk}\hat g_1(\kk_1)\hat g_2(\kk_2).
\label{eq:conv}
\end{equation}
For a statistically homogeneous random field, $\langle\hat g(\kk)\hat g^*(\kk')\rangle=L^3\,P_g(k)\,\delta_{\kk\kk'}$ defines the power spectrum, and for two fields $P_{ab}$ is defined the same way; the cross-correlation coefficient is $r_{ab}(k)=P_{ab}/\sqrt{P_{aa}P_{bb}}$.

The Nyquist wavenumber of level $\ell$ is $k_{\Ny,\ell}=\pi/h_\ell$. The space of functions band-limited to level $\ell$ is
\begin{equation}
\mathcal B_\ell=\{g:\ \hat g(\kk)=0\ \text{unless}\ |k_j|<k_{\Ny,\ell}\ \text{for } j=1,2,3\}.
\end{equation}
(Strict inequality: the Nyquist planes are excluded. In the code they are zeroed explicitly; this loses a set of measure zero of modes and removes the usual ambiguity of the $\pm k_\Ny$ identification.)

\subsection{Sampling a band-limited function on an offset grid}

Let $g\in\mathcal B_\ell$ and define its discrete Fourier transform from the samples on $\Lambda_\ell$ exactly as \texttt{numpy.fft} does, i.e.\ with phases referenced to the integer index and not to the physical position:
\begin{equation}
G_\ell(\kk):=\sum_{\mathbf i}g(\q_{\mathbf i})\,e^{-i\kk\cdot\mathbf i h_\ell},\qquad \kk\in k_F\Z^3,\ |k_j|<k_{\Ny,\ell}.
\end{equation}

\begin{lemma}[DFT of a band-limited function]\label{lem:dft}
For $g\in\mathcal B_\ell$ and $|k_j|<k_{\Ny,\ell}$,
\begin{equation}
G_\ell(\kk)=h_\ell^{-3}\,e^{\,i\kk\cdot o\mathbf 1 h_\ell}\,\hat g(\kk).
\label{eq:dft}
\end{equation}
\end{lemma}
\begin{proof}
Insert \eqref{eq:fourier} into the definition:
$G_\ell(\kk)=L^{-3}\sum_{\kk'}\hat g(\kk')e^{i\kk'\cdot o\mathbf 1h_\ell}\sum_{\mathbf i}e^{i(\kk'-\kk)\cdot\mathbf i h_\ell}$. The inner sum equals $N_\ell^3$ if $\kk'-\kk\in 2k_{\Ny,\ell}\Z^3$ and $0$ otherwise. Since both $\kk$ and $\kk'$ lie strictly inside the Nyquist cube, the only solution is $\kk'=\kk$. Then $G_\ell(\kk)=N_\ell^3L^{-3}e^{i\kk\cdot o\mathbf 1h_\ell}\hat g(\kk)$ and $N_\ell^3/L^3=h_\ell^{-3}$.
\end{proof}

\begin{corollary}[Coarse and fine DFTs of the same function]\label{cor:cf}
If $g\in\mathcal B_c$ (hence also $g\in\Bf$), then for $\kk$ inside the coarse Nyquist cube
\begin{equation}
G_c(\kk)=\Big(\frac{h_f}{h_c}\Big)^{3}e^{\,i\kk\cdot o\mathbf 1(h_c-h_f)}\,G_f(\kk)=\frac18\,e^{\,i\kk\cdot o\mathbf 1 h_f}\,G_f(\kk).
\label{eq:cfphase}
\end{equation}
\end{corollary}
This is the factor $(N_c/N_f)^3$ and the phase $e^{+i\kk\cdot o\mathbf 1h_f}$ used in the function \texttt{restrict\_spectral}; \texttt{prolong\_spectral} uses the inverse. With $o=0$ the coarse points are a subset of the fine points and the phase is $1$; with $o=\tfrac12$ the coarse points sit at the centres of $2\times2\times2$ blocks of fine points and the half-fine-cell phase is what makes the two grids consistent.

\begin{exercise}
Verify \eqref{eq:cfphase} directly for the single mode $g(\q)=e^{i\kk\cdot\q}$ with $o=\tfrac12$ by writing $\q^{c}_{\mathbf i}=(2\mathbf i+\mathbf 1)h_f$.
\end{exercise}

%======================================================================
\section{Restriction, prolongation and the coarse/detail decomposition}
%======================================================================

\subsection{Spectral projectors}

Let $W:k_F\Z^3\to\{0,1\}$ be an indicator supported strictly inside the coarse Nyquist cube. The two choices used in the code are
\begin{equation}
W_{\rm cube}(\kk)=\mathbb 1[\,|k_j|<k_{\Ny,c}\ \forall j\,],\qquad
W_{\rm sph}(\kk)=\mathbb 1[\,|\kk|<\alpha\,k_{\Ny,c}\,]\,W_{\rm cube}(\kk),\quad \alpha\le1 .
\end{equation}
Define the Fourier multiplier $\PW$ by $\widehat{\PW g}=W\hat g$ and the subspace $\BW=\PW\Bf=\{g\in\Bf:\ \hat g=W\hat g\}$.

\begin{proposition}\label{prop:proj}
$\PW$ is an orthogonal projector on $\Bf$ (with the $L^2$ inner product), $\PW^2=\PW$, $\PW^\dagger=\PW$, and
\begin{equation}
\Bf=\BW\oplus\BW^{\perp},\qquad \BW^\perp=(I-\PW)\Bf .
\end{equation}
Every $g\in\Bf$ has the unique orthogonal decomposition $g=\PW g+(I-\PW)g$.
\end{proposition}
\begin{proof}
$W^2=W$ because $W\in\{0,1\}$. Self-adjointness is Parseval's identity:
\begin{equation*}
\langle g_1,\PW g_2\rangle=L^{-3}\sum_\kk \hat g_1^*\,W\,\hat g_2=\langle \PW g_1,g_2\rangle .
\end{equation*}
\end{proof}

\subsection{$R$ and $P$}

\begin{definition}[Restriction and prolongation]
Identify a coarse array with the unique function of $\Bc$ that interpolates it (Lemma~\ref{lem:dft} makes this identification explicit, including the offset phase). Define
\begin{align}
R&:\Bf\to\BW,\qquad R g=\PW g \ \text{(represented by its samples on }\Lambda_c),\\
P&:\BW\to\Bf,\qquad P g=g \ \text{(represented by its samples on }\Lambda_f).
\end{align}
\end{definition}
In words: $R$ keeps the modes selected by $W$ and re-samples on the coarse grid; $P$ is the band-limited (``zero-padding'') interpolation from the coarse grid to the fine grid. The array operations are exactly the coefficient copy with the factor and phase of Corollary~\ref{cor:cf}.

\begin{proposition}[Exact identities]\label{prop:RP}
$RP=I$ on $\BW$, and $PR=\PW$ on $\Bf$.
\end{proposition}
\begin{proof}
For $g\in\BW$, $Pg=g$ and $R g=\PW g=g$. For $g\in\Bf$, $PRg=\PW g$.
\end{proof}

\begin{remark}[The coarse state of the chain; cube versus sphere]\label{rem:clean}
The natural state variable of the progressive chain at a level is an element of $\BW$. With $W=W_{\rm cube}$ this is simply the coarse array: nothing is discarded, the nested-initial-condition split of Section~\ref{sec:linear} coincides with the model split, and every statement below about ``the new randomness'' is exact. This is the default.

The sphere is tempting because it is isotropic and because the corner modes of the coarse cube ($|\kk|>\alpha k_{\Ny,c}$; for $\alpha=1$ they are $1-\pi/6\approx48\%$ of the coarse modes) are the ones a low-resolution run resolves worst. Using it, however, means projecting the input, $\Psic\leftarrow\PW\Psic$, and re-generating the corners as if they were fresh randomness. That is an approximation: the corner modes $\chi$ of the coarse initial conditions are independent of the rest \emph{a priori}, but the nonlinear coarse output inside the sphere, $\PW\Nbody_c[\delta_W+\chi]$, carries information about $\chi$ through mode coupling, so $p(\chi\mid\PW\Psic)\neq p(\chi)$. The exact treatment keeps the corners in the state (cube) and lets the correction band repair them. The sphere should be regarded as an ablation, not as the theoretical baseline.
\end{remark}

\begin{remark}[Why $W$ must be an indicator]
If $W$ is replaced by a smooth taper $0\le W\le1$, then $\PW$ is no longer a projector ($W^2\neq W$), $RP\neq I$, and the ``coarse'' and ``detail'' parts overlap in the taper band. A taper can still be used as a partition of unity in the \emph{source distribution} of the generative model (Section~\ref{sec:fm}), but not to define the state of the chain.
\end{remark}

\subsection{The decomposition of a fine field and the two tasks of a $2\times$ step}

For the fine field $\Psif\in\Bf$ (each component), Proposition~\ref{prop:proj} gives
\begin{equation}
\boxed{\ \Psif=P\,(R\Psif)\;+\;d,\qquad d:=(I-\PW)\Psif\in\BW^\perp .\ }
\label{eq:decomp}
\end{equation}
The first term is what a coarse grid can represent; $d$ is the \emph{detail}. Given a coarse simulation $\Psic$ (cleaned as in Remark~\ref{rem:clean}), define the \emph{coarse correction}
\begin{equation}
\varepsilon:=R\Psif-\Psic .
\label{eq:eps}
\end{equation}
Then $\Psif=P(\Psic+\varepsilon)+d$: a $2\times$ step must (i) produce $\varepsilon$, the error of the coarse run on the scales it resolves, and (ii) produce $d$. Both are functions of the same fine-grid residual $\Psif-P\Psic$: $\varepsilon$ is its $W$-band, $d$ its complementary band. The model therefore has a single fine-grid output; the two ``heads'' are Fourier bands of it.

\subsection{Counting the new degrees of freedom}

Per component, $\dim\Bf-\dim\BW$ is the number of modes outside $W$ but inside the fine cube. For $W_{\rm cube}$ and ignoring the Nyquist planes, $(2N_c)^3-N_c^3=7N_c^3$, which is also the number of fine grid points not shared with a coarse cell centre ($8-1$ per coarse cell). The new \emph{randomness} of a $2\times$ step is a scalar field with this many modes (the octave of the initial conditions, Section~\ref{sec:push}); the detail $d$ itself has three components and, as Section~\ref{sec:2lpt} shows, also contains a deterministic part generated by the coarse modes alone.

\subsection{Haar restriction for comparison}

The Haar (block-mean) restriction $R_H\Psi(\q^c_{\mathbf i})=\tfrac18\sum_{\mathbf s\in\{0,1\}^3}\Psi(\q^f_{2\mathbf i+\mathbf s})$ conserves mass and centre of mass exactly but is not a spectral projector.

\begin{proposition}[Transfer function and aliasing of Haar]\label{prop:haar}
Let $o=\tfrac12$ (block centres coincide with coarse points). For a single mode $e^{i\kk\cdot\q}$,
\begin{equation}
R_H e^{i\kk\cdot\q}\big|_{\q^c}=T_H(\kk)\,e^{i\kk\cdot\q^c},\qquad T_H(\kk)=\prod_{j=1}^3\cos\!\Big(\frac{k_jh_f}{2}\Big).
\end{equation}
Consequently, for a band-limited fine field, $R_H\Psif=\mathrm{Samp}_c\big[T_H\hat\Psif\big]$ where the sampling aliases every mode outside the coarse cube onto a mode inside it with weight $T_H\neq0$.
\end{proposition}
\begin{proof}
The fine points of a block are $\q^c+(\mathbf s-\tfrac12\mathbf 1)h_f$; averaging $e^{ik_j(s_j-1/2)h_f}$ over $s_j\in\{0,1\}$ gives $\cos(k_jh_f/2)$ per axis.
\end{proof}

At first order in perturbation theory (Section~\ref{sec:linear}) this means $R_H\Psi^{(1)}_f-\Psi^{(1)}_c=(T_H-1)\Psi^{(1)}_c+\text{(aliased octave modes)}$: a deterministic mismatch $\approx-\tfrac18\sum_j(k_jh_f)^2\Psi_c^{(1)}$ at low $k$, with power $\propto k^4P_{\Psi}(k)$, plus a \emph{random} term that the spectral $R$ does not have. (With $o=0$ the block centre is displaced by $\tfrac12h_f\mathbf 1$ from the coarse point, an extra phase $e^{i\kk\cdot\mathbf 1h_f/2}$ appears, and the mismatch is $O(k)$ with power $\propto k^2P_\Psi$.) These are the flat/negative low-$k$ slopes of the Haar curves in the self-test.

\begin{exercise}
For $P_\Psi\propto k^{n-2}$ with $n=-1.5$, predict the low-$k$ slope of $P_\varepsilon$ for Haar with $o=\tfrac12$ and with $o=0$, and compare with the measured $\approx+0.6$ and $\approx-1.7$.
\end{exercise}

%======================================================================
\section{Nested initial conditions and linear theory}\label{sec:linear}
%======================================================================

\subsection{Nested initial conditions}

The linear density contrast is a Gaussian random field with $\langle\hat\delta(\kk)\hat\delta^*(\kk')\rangle=L^3P_\lin(k)\delta_{\kk\kk'}$ (normalised to $z=0$; the growth factor $D(a)$ multiplies it at other times). The initial-condition generator of level $\ell$ draws exactly the modes of $\mathcal B_\ell$, using the same random numbers for the same $\kk$ at every level. Hence
\begin{equation}
\delta_f=\delta_c+\eta,\qquad \hat\delta_c=W_{\rm cube}\hat\delta_f,\qquad \hat\eta=(1-W_{\rm cube})\hat\delta_f ,
\label{eq:nested}
\end{equation}
and $\eta$, the \emph{octave}, is independent of $\delta_c$ (disjoint sets of independent Gaussian modes).

\subsection{Zel'dovich displacement}

At first order, $\Psi^{(1)}=-D\nabla\phi^{(1)}$ with $\nabla^2\phi^{(1)}=\delta$, i.e.
\begin{equation}
\hat\Psi^{(1)}(\kk)=D\,\frac{i\kk}{k^2}\hat\delta(\kk)=:D\,\Zel\hat\delta ,\qquad \nabla\cdot\Psi^{(1)}=-D\delta .
\label{eq:zel}
\end{equation}
$\Zel$ is a Fourier multiplier; it commutes with every spectral projector.

\begin{proposition}[Linear theory of the correction and the detail]\label{prop:linear}
With nested initial conditions and any indicator $W\le W_{\rm cube}$:
\begin{enumerate}
\item[(i)] $\varepsilon^{(1)}:=R\Psi^{(1)}_f-\PW\Psi^{(1)}_c=0$ identically.
\item[(ii)] $d^{(1)}=(I-\PW)\Psi_f^{(1)}=D\Zel(1-W)\hat\delta_f$ is a Gaussian field, independent of $\PW\Psi_c^{(1)}$, with covariance diagonal in $\kk$ and equal to $D^2P_\lin(k)/k^2$ on the modes outside $W$.
\item[(iii)] The velocity obeys the same statements with $D\to \dot D=aHfD$.
\end{enumerate}
\end{proposition}
\begin{proof}
(i) $R\Psi^{(1)}_f=\PW D\Zel\hat\delta_f=D\Zel W\hat\delta_f$ and $\PW\Psi^{(1)}_c=D\Zel WW_{\rm cube}\hat\delta_f=D\Zel W\hat\delta_f$ because $W\le W_{\rm cube}$. (ii) $(1-W)\hat\delta_f$ and $W\hat\delta_f$ involve disjoint sets of independent Gaussian coefficients. (iii) $\mathbf v^{(1)}=\dot D\Zel\hat\delta$.
\end{proof}

\begin{remark}
(i) is the reason to prefer a spectral $R$ over Haar: with the spectral $R$ the whole correction $\varepsilon$ is a mode-coupling effect (Section~\ref{sec:2lpt}), while Haar creates a first-order mismatch. (ii) is the reason to hard-code the linear octave into the model: at early times or on large scales the detail is \emph{exactly} a known Gaussian, and the network should only have to learn the departure from it.
\end{remark}

%======================================================================
\section{Second order: the structure of the conditional}\label{sec:2lpt}
%======================================================================

\subsection{2LPT}

In the conventions of Bouchet et al.\ (1995),
\begin{equation}
\Psi=-D\nabla\phi^{(1)}+D_2\nabla\phi^{(2)},\qquad \nabla^2\phi^{(2)}=S:=\sum_{i>j}\big[\phi^{(1)}_{,ii}\phi^{(1)}_{,jj}-(\phi^{(1)}_{,ij})^2\big],\qquad D_2\simeq-\tfrac37D^2 .
\label{eq:2lpt}
\end{equation}
Two identities will be used. First, with $D=1$ and the deformation tensor $D_{ij}:=\partial_i\Psi_j$, whose first-order value is $D^{(1)}_{ij}=-\phi^{(1)}_{,ij}$,
\begin{equation}
S=\tfrac12\Big[(\nabla^2\phi^{(1)})^2-\phi^{(1)}_{,ij}\phi^{(1)}_{,ij}\Big]=\tfrac12\Big[(\tr D^{(1)})^2-D^{(1)}_{ij}D^{(1)}_{ij}\Big]=\sum_{i<j}\lambda_i\lambda_j ,
\label{eq:S}
\end{equation}
where $\lambda_i$ are the eigenvalues of $D^{(1)}$ (the second elementary symmetric polynomial). Second, in Fourier space, using $\hat\phi^{(1)}=-\hat\delta/k^2$ and \eqref{eq:conv},
\begin{equation}
\hat S(\kk)=\frac{1}{2L^3}\sum_{\kk_1+\kk_2=\kk}K(\kk_1,\kk_2)\,\hat\delta(\kk_1)\hat\delta(\kk_2),\qquad
K(\kk_1,\kk_2)=1-\frac{(\kk_1\cdot\kk_2)^2}{k_1^2k_2^2}=\frac{|\kk_1\times\kk_2|^2}{k_1^2k_2^2},
\label{eq:kernel}
\end{equation}
and the second-order displacement is the gradient of the inverse Laplacian of $S$. Since the Fourier multiplier of $\nabla^{-2}$ is $-1/k^2$, the real-space and Fourier-space forms are (with $D=1$)
\begin{equation}
\Psi^{(2)}=D_2\nabla\nabla^{-2}S=-\tfrac37\,\nabla\nabla^{-2}S,\qquad
\hat\Psi^{(2)}(\kk)=D_2\,i\kk\,\hat\phi^{(2)}(\kk)=-D_2\frac{i\kk}{k^2}\hat S(\kk)=+\frac37\,\frac{i\kk}{k^2}\hat S(\kk).
\label{eq:psi2}
\end{equation}
(The two signs differ because $\nabla^{-2}\to-1/k^2$; an earlier version of these notes dropped the minus sign in the real-space form.) We write $\mathcal T_2:=D_2\nabla\nabla^{-2}$ for the linear operator that turns a source into a second-order displacement.

\begin{exercise}
Derive \eqref{eq:S} by expanding $(\sum_i\phi_{,ii})^2$ and $\phi_{,ij}\phi_{,ij}$; derive \eqref{eq:kernel} from $\phi_{,ij}\to k_ik_j\hat\delta/k^2$.
\end{exercise}

\subsection{The long/short split}

Write $\delta_f=\delta_L+\eta$ with $\delta_L:=\delta_c$ (the coarse, or ``long'', modes). Since $S$ is a quadratic form, its polarisation gives
\begin{equation}
S[\delta_L+\eta]=S[\delta_L]+S[\eta]+B(\delta_L,\eta),\qquad
B(a,b):=\tr D^{a}\,\tr D^{b}-D^{a}_{ij}D^{b}_{ij},
\label{eq:split}
\end{equation}
with $D^a_{ij}=-\partial_i\partial_j\nabla^{-2}a$. The coarse simulation, if it solved the continuum equations exactly for the band-limited input $\delta_L$, would produce $\Psi^{(2)}_c=\mathcal T_2S[\delta_L]$, while the fine one produces $\Psi^{(2)}_f=\mathcal T_2\big[S[\delta_L]+B(\delta_L,\eta)+S[\eta]\big]$ (all outputs understood as band-limited to the fine cube). Projecting with $\PW$ and $I-\PW$:
\begin{equation}
\boxed{\ \varepsilon^{(2)}=\PW\,\mathcal T_2\big[B(\delta_L,\eta)+S[\eta]\big],\qquad
d^{(2)}=(I-\PW)\,\mathcal T_2\big[S[\delta_L]+B(\delta_L,\eta)+S[\eta]\big]\ }
\label{eq:eps2}
\end{equation}
The $S[\delta_L]$ term cancels in the \emph{correction} (both runs contain it on the coarse band) but \emph{not} in the detail: two coarse modes can add up to a mode outside $W$. A concrete instance, in units of $k_F$ with coarse cutoff $4$: $\mathbf p=(3,2,0)$ and $\mathbf q=(3,-2,0)$ are both inside the coarse cube, $\mathbf p+\mathbf q=(6,0,0)$ is outside it but inside the fine cube, and $K(\mathbf p,\mathbf q)=1-25/169=144/169\neq0$. So even with $\eta\equiv0$ the fine detail is non-zero. Consequences: (i) the detail band contains a part that is a deterministic function of the coarse initial conditions (the nonlinear ``harmonics'' of the coarse modes), on top of the part driven by the octave; the model, being conditioned on the coarse field, produces both, but the sentence ``the detail is the new randomness'' is true only at first order; (ii) $d^{(2)}$ is not band-limited to $2\times$ the coarse cutoff in any simple way, which is one more reason to give the network the whole fine grid rather than a per-cell rule.

\subsection{Suppression of the correction at low $k$}

\begin{lemma}[Exact small-$k$ form of the kernel]\label{lem:kernel}
For $\kk_2=\kk-\kk_1$ and $\mu:=\hat\kk\cdot\hat\kk_1$,
\begin{equation}
K(\kk_1,\kk-\kk_1)=\frac{k^2\,(1-\mu^2)}{|\kk-\kk_1|^2}.
\end{equation}
\end{lemma}
\begin{proof}
$\kk_1\times\kk_2=\kk_1\times(\kk-\kk_1)=\kk_1\times\kk$, so $|\kk_1\times\kk_2|^2=k_1^2k^2(1-\mu^2)$; divide by $k_1^2k_2^2$.
\end{proof}
(An equivalent route: with $\epsilon=k/k_1$, $(\kk_1\cdot\kk_2)^2/(k_1^2k_2^2)=(1-\epsilon\mu)^2/(1-2\epsilon\mu+\epsilon^2)=1-\epsilon^2(1-\mu^2)/(1-2\epsilon\mu+\epsilon^2)$.)

\begin{proposition}[The $k^2$/$k^4$ law for the octave--octave term]\label{prop:k2}
Let $\vartheta^{(2)}:=-\nabla\cdot\Psi^{(2)}$ denote the divergence of the second-order displacement (it is \emph{not} the full second-order Eulerian density, see Remark~\ref{rem:density}). For the octave--octave source, Lemma~\ref{lem:kernel} gives, for every realisation,
\begin{equation}
\hat S_{\eta\eta}(\kk)=k^2\,\hat Q(\kk),\qquad \hat Q(\kk)=\frac{1}{2L^3}\sum_{\kk_1}\frac{1-\mu_1^2}{|\kk-\kk_1|^2}\,\hat\eta(\kk_1)\hat\eta(\kk-\kk_1),
\end{equation}
where both $\kk_1$ and $\kk-\kk_1$ lie outside the coarse cube, so the denominators are bounded below by $k_{\Ny,c}^2$. The statistical statement follows by Wick's theorem (two contractions, symmetric kernel):
\begin{align}
P_{S_{\eta\eta}}(k)&=\frac12\int\frac{\dd^3p}{(2\pi)^3}\,K^2(\mathbf p,\kk-\mathbf p)\,P_\eta(p)\,P_\eta(|\kk-\mathbf p|)\nonumber\\
&=k^4\cdot\frac12\int\frac{\dd^3p}{(2\pi)^3}\frac{(1-\mu_p^2)^2}{|\kk-\mathbf p|^4}P_\eta(p)P_\eta(|\kk-\mathbf p|)
\ \xrightarrow[k\to0]{}\ C\,k^4 ,
\end{align}
with $0<C<\infty$ because the remaining integral is over a band away from the origin and converges to a constant. Hence, using \eqref{eq:psi2},
\begin{equation}
P_{\Psi^{(2)}_{\eta\eta}}(k)=\Big(\tfrac37\Big)^2\frac{P_{S_{\eta\eta}}(k)}{k^2}\propto k^2,\qquad
P_{\vartheta^{(2)}_{\eta\eta}}(k)=\Big(\tfrac37\Big)^2P_{S_{\eta\eta}}(k)\propto k^4 .
\end{equation}
\end{proposition}
The realisation-level statement is that $\hat Q(\kk)$ is a random variable whose variance is $k$-independent for $k\ll k_{\Ny,c}$; the ensemble statement above is the precise form, and it is what the shell-averaged spectra in the code measure.

Three qualifications. First, the mixed term $B(\delta_L,\eta)$ carries the same kernel factor, but the pairs $(\kk_L,\kk_\eta)$ that can add up to a small $\kk$ must both lie within a distance $\sim k$ of the cutoff surface, so its phase space shrinks with $k$ as well and its low-$k$ power falls at least as fast as the octave--octave term; the low-$k$ correction is dominated by $S[\eta]$. Second, this is a statement about the ideal, de-aliased, second-order contribution; higher-order deterministic responses and the discreteness errors of a real coarse run (Remark~\ref{rem:disc}) need not share the exponent. Third, the divergence $\vartheta$ and the Eulerian density are different fields; see the next remark.

\begin{remark}[Divergence versus Eulerian density]\label{rem:density}
Mass conservation gives, for $\kk\neq0$, $\hat\delta_E(\kk)=\int\dd^3q\,e^{-i\kk\cdot\q}\big[e^{-i\kk\cdot\Psi(\q)}-1\big]$, whose second-order part is
$\hat\delta^{(2)}_E=-ik_i\hat\Psi^{(2)}_i-\tfrac12k_ik_j\,\widehat{\Psi^{(1)}_i\Psi^{(1)}_j}$. The second term (first-order displacements changing the spacing of mass elements) is absent from $\vartheta^{(2)}$. Assembling both gives the standard $F_2$ kernel of Eulerian perturbation theory, for which $F_2(\mathbf p,\kk-\mathbf p)=O(k^2)$ as $\kk\to0$ with $\mathbf p$ fixed; so the octave--octave part of the true density power also falls as $k^4$ (Peebles' tail), but that is a separate calculation. In the code the ``divergence'' spectra are exactly $\vartheta$; density spectra from particle deposition are a different diagnostic.
\end{remark}

This is the Lagrangian form of mass and momentum conservation: rearranging matter inside a small region cannot move mass on large scales except at order $k^2$ in the density. In the language of the effective field theory of large-scale structure, $\frac37\nabla Q$ is the stochastic term; the ``sound-speed'' counterterm $\propto c_s^2k^2\hat\Psi^{(1)}$ appears one order higher, from $\delta_L\eta\eta$. The de-aliased 2LPT self-test measures shell-averaged slopes $2.19$ (displacement) and $4.14$ (divergence).

\begin{remark}[Discreteness caveat]\label{rem:disc}
The cancellation of $S[\delta_L]$ in the correction of \eqref{eq:eps2} assumes the coarse run computes the quadratic term without aliasing. A grid code evaluates products on its own grid, and a product of two modes near $k_{\Ny,c}$ has support up to $2k_{\Ny,c}$, which aliases back to low $k$ with the kernel evaluated at the \emph{large} arguments, i.e.\ without the $k^2$ suppression. A real $N$-body run has its own version of this (PM aliasing, softening, time stepping). Empirically it is a small white floor at low $k$; its amplitude is what \texttt{phase0\_octaves.py} panel (a) measures on the real $64/128/256/512$ series.
\end{remark}

\subsection{How the octave couples to the coarse field: locality and the deformation tensor}

In real space the cross term of \eqref{eq:split} is pointwise:
\begin{equation}
B(\delta_L,\eta)(\q)=\tr D^{L}(\q)\,\tr D^{\eta}(\q)-D^{L}_{ij}(\q)\,D^{\eta}_{ij}(\q),
\label{eq:B}
\end{equation}
where $D^L_{ij}=\partial_i\Psi^{(1)}_{L,j}$ is the (first-order) deformation tensor of the coarse field and $D^\eta$ that of the linear octave. Thus, to leading order, the coarse field enters the octave-driven part of the detail \emph{only} through its deformation tensor, locally in $\q$, followed by the potential solve $\mathcal T_2$. (The coarse-only harmonics $(I-\PW)\mathcal T_2S[\delta_L]$ of \eqref{eq:eps2} are likewise a local quadratic function of $D^L$ followed by $\mathcal T_2$.) Decompose $D^L$ into its trace and traceless parts, $D^L_{ij}=-\tfrac13\delta_L\,\delta_{ij}+\tau^L_{ij}$ (using $\tr D^L=-\delta_L$):
\begin{equation}
B=-\tfrac23\,\delta_L\,\tr D^\eta-\tau^L_{ij}D^\eta_{ij}.
\end{equation}
The first term is the isotropic (``separate universe'') coupling: a local overdensity enhances the growth of the small-scale modes; the second is the tidal coupling, which is what anisotropic separate-universe simulations measure. If $D^L$ is constant over a patch, the short modes in that patch evolve as in a homogeneous, anisotropically deformed background.

This is the justification for the network inputs: the coarse field is presented as $D_{ij}(P\Psic)$ (nine dimensionless channels), never as $\Psic$ itself. A constant displacement of the coarse field (a pure translation) changes nothing, and the leading physics is a local function of $D_{ij}$.

\begin{exercise}
Starting from \eqref{eq:B}, show that the Fourier-space kernel of the cross term is $K(\kk_L,\kk_\eta)=1-(\hat\kk_L\cdot\hat\kk_\eta)^2$ and interpret $(\hat\kk_L\cdot\hat\kk_\eta)^2=\hat k_{\eta,i}\hat k_{\eta,j}\hat k_{L,i}\hat k_{L,j}$ as the contraction of the octave direction with the coarse deformation tensor.
\end{exercise}

%======================================================================
\section{The conditional distribution: exact statements}\label{sec:push}
%======================================================================

Let $\Nbody_\ell$ denote the map of the level-$\ell$ simulation code from its linear initial conditions $\delta\in\mathcal B_\ell$ to the pair $(\Psi_\ell,\mathbf v_\ell)$ at the output time. The simulation is deterministic.

\begin{proposition}[Push-forward form of the physical conditional]\label{prop:push}
Let $C$ denote the coarse output that is actually conditioned on. If $C$ determines $\delta_c$ (for instance $C=(\Psic,\mathbf v_c)$, the full phase-space output, and $\Nbody_c$ is injective), then
\begin{equation}
\Psif\ \overset{d}{=}\ \Nbody_f[\delta_c+\eta],\qquad \eta\sim\mathcal N\big(0,\ P_\lin\ \text{restricted to the octave}\big),\quad \eta\perp\delta_c ,
\end{equation}
i.e.\ $p(\Psif\mid C)$ is the push-forward of a known Gaussian with $7N_c^3$ scalar degrees of freedom through the map $\eta\mapsto\Nbody_f[\delta_c+\eta]$. In general, without assuming that $C$ determines $\delta_c$,
\begin{equation}
p(\Psif\in A\mid C)=\int\!\!\int\mathbb 1_A\big(\Nbody_f[\delta_c+\eta]\big)\,\gamma(\dd\eta)\,p(\dd\delta_c\mid C),
\label{eq:mixture}
\end{equation}
a mixture over the residual uncertainty $p(\delta_c\mid C)$ of the coarse initial conditions given $C$, with $\gamma$ the octave Gaussian.
\end{proposition}
\begin{proof}
By \eqref{eq:nested}, $\delta_f=\delta_c+\eta$ with $\eta$ independent of $\delta_c$, hence of any function $C$ of $\delta_c$; condition on $C$ and integrate over what it leaves undetermined.
\end{proof}

\begin{remark}[What is actually conditioned on]
The state used by the model is not the full coarse output. It is successively reduced: $(\Psic,\mathbf v_c)\to\Psic\to\PW\Psic\to$ a local patch of $D_{ij}(P\Psic)$ seen by a finite receptive field. Injectivity of the full map $\Nbody_c$ says nothing about whether these reduced observations determine $\delta_c$; the honest description is \eqref{eq:mixture}, and $p(\delta_c\mid C)$ is degenerate only in the idealised case. Two practical consequences. (i) The randomness a $2\times$ step must supply is the octave \emph{plus} whatever the reduced conditioning leaves unresolved of $\delta_c$; only the former has the clean Gaussian form. (ii) With the sphere window the discarded corner modes $\chi$ are part of $\delta_c$; the retained field $\PW\Nbody_c[\delta_W+\chi]$ is informative about $\chi$ through mode coupling, so $\chi$ is \emph{not} an independent Gaussian given $C$, and a source that re-samples the corners from the prior is an approximation (Remark~\ref{rem:clean}). With the cube window this issue is absent.
\end{remark}

\begin{corollary}[Markov property across levels]
For the full nested chain, $p(\Psi_{\ell+2}\mid\Psi_{\ell+1},\Psi_\ell)=p(\Psi_{\ell+2}\mid\Psi_{\ell+1})$ whenever $\Psi_{\ell+1}$ determines $\delta_{\ell+1}$, so the autoregressive chain is exact:
$p(\Psi_{L}\mid\Psi_0)=\int\prod_{\ell=0}^{L-1}p(\Psi_{\ell+1}\mid\Psi_\ell)\,\dd\Psi_1\cdots\dd\Psi_{L-1}$.
\end{corollary}
\begin{proof}
$\delta_{\ell+1}\supset\delta_\ell$, and $\Psi_{\ell+2}$ depends on $\delta_{\ell+1}$ and an independent octave.
\end{proof}
The Markov property is a property of the full (hidden) state. A chain built on reduced states (projected fields, patches) is Markov only if the reduction is a sufficient statistic; otherwise earlier levels carry information the current state has lost, and the one-step conditional used in the chain is a modelling approximation. This is one reason to keep the cube window and, eventually, to include the velocity field in the state.

Two consequences shape the method. First, the randomness of a $2\times$ step has a physical identity (the octave initial conditions) and a known dimension; a model whose noise is that octave can be trained \emph{with the true $\eta$} (the setting of Paper IV, extended level by level) and used with a sampled $\eta$ (a generative model). Second, under the idealised conditioning of Proposition~\ref{prop:push}, $\Psif$ is a deterministic function of $(C,\eta)$: a perfect regressor with the true $\eta$ as input would be exact, and fed a fresh $\eta$ it would be a valid generative model. Chaos in multi-stream regions does not change this in exact arithmetic (a fixed initial condition has one outcome); what it changes is how much of $(\delta_c,\eta)$ a model with finite capacity, finite receptive field and finite training data can effectively use, and how fast round-off is amplified in a real code. The conditional flow matching of the next section is formulated so that its correctness in generative mode does not rest on the regressor being exact.

%======================================================================
\section{Conditional flow matching with the physical coupling}\label{sec:fm}
%======================================================================

Everything in this section is conditional on the coarse field $C$ (the coarse array with the cube window; the projected field if the sphere is used); the conditioning is suppressed in the notation except where it matters. All fields are on the fine grid and measured in units of $h_f$.

\subsection{Source, target, coupling}

\begin{align}
x_0&=P\Psic+\Psi^{(1)}[\eta]=P\Psic+D\Zel\hat\eta,\qquad \eta\sim\mathcal N(0,P_\lin|_{\rm octave}),\label{eq:x0}\\
x_1&=\Psif .
\end{align}
Let $q_0,q_1$ be the laws of $x_0,x_1$. A \emph{coupling} is a joint law $\pi(x_0,x_1)$ with these marginals. Two couplings matter:
\begin{itemize}
\item[] \emph{independent}: $\eta$ in $x_0$ is drawn independently of the $\eta$ that produced $x_1$;
\item[] \emph{physical}: the same $\eta$ appears in both, $x_1=\Nbody_f[\delta_c+\eta]$ and $x_0=P\Psic+D\Zel\hat\eta$.
\end{itemize}
Both have the correct marginals, because the marginal of $x_0$ only depends on the law of $\eta$, which is the same.

\subsection{The interpolant and the marginal velocity}

Define the straight-line interpolant and its velocity
\begin{equation}
x_t=(1-t)x_0+t\,x_1,\qquad \dot x_t=x_1-x_0,\qquad t\in[0,1],
\end{equation}
let $p_t$ be the law of $x_t$ under $\pi$ (so $p_0=q_0$, $p_1=q_1$), and define the \emph{marginal velocity field}
\begin{equation}
u_t(x):=\E_\pi\big[x_1-x_0\ \big|\ x_t=x\big].
\label{eq:ut}
\end{equation}

\begin{theorem}[Marginals are transported by $u_t$, for any coupling]\label{thm:fm}
$p_t$ solves the continuity equation $\partial_tp_t+\nabla\cdot(p_tu_t)=0$. Hence, if $u_t$ is Lipschitz, the flow $\dot X_t=u_t(X_t)$, $X_0\sim q_0$, satisfies $X_t\sim p_t$ and in particular $X_1\sim q_1$. The result depends on $\pi$ only through $u_t$; the marginals of the generated samples do not depend on the coupling.
\end{theorem}
\begin{proof}
For any smooth compactly supported test function $\varphi$,
\begin{equation}
\frac{\dd}{\dd t}\E\,\varphi(x_t)=\E\big[\nabla\varphi(x_t)\cdot(x_1-x_0)\big]
=\E\Big[\nabla\varphi(x_t)\cdot\E[x_1-x_0\mid x_t]\Big]=\int\nabla\varphi(x)\cdot u_t(x)\,p_t(x)\,\dd x,
\end{equation}
using the tower property of conditional expectation. The left side is $\frac{\dd}{\dd t}\int\varphi\,p_t$, so $\int\varphi\,\partial_tp_t=-\int\varphi\,\nabla\cdot(p_tu_t)$ for all $\varphi$, which is the weak continuity equation. The law of $X_t$ for the ODE flow solves the same equation with the same initial condition; uniqueness for Lipschitz vector fields gives $\mathrm{Law}(X_t)=p_t$.
\end{proof}

\subsection{The training objective}

\begin{proposition}[Minimiser of the flow-matching loss]\label{prop:loss}
For $v_\theta:\R^{d}\times[0,1]\to\R^d$ and $t\sim\mathcal U[0,1]$,
\begin{equation}
\mathcal L(\theta)=\E_{t,\pi}\big\|v_\theta(x_t,t)-(x_1-x_0)\big\|^2
=\E_{t}\,\E_{x_t\sim p_t}\big\|v_\theta(x_t,t)-u_t(x_t)\big\|^2+\E_t\,\E\big[\Var(x_1-x_0\mid x_t)\big].
\end{equation}
The second term does not depend on $\theta$; the minimiser is $v_\theta=u_t$.
\end{proposition}
\begin{proof}
Bias--variance decomposition of the squared error around the conditional mean $u_t(x_t)=\E[x_1-x_0\mid x_t]$; the cross term vanishes by the tower property.
\end{proof}

In the code, $v_\theta$ receives $x_t-P\Psic$ (not $x_t$), the nine channels of $D_{ij}(P\Psic)$, and $t$; the target is $x_1-x_0$; everything is in units of $h_f$. Since $x_1-x_0$ is unchanged by a constant shift of all fields and $D_{ij}$ is shift invariant, the learned map is translation invariant by construction.

\subsection{Why the physical coupling is the right one}

By Theorem~\ref{thm:fm} the physical coupling is admissible. It is preferable because it makes the transport short and nearly straight. Write $\Psif=P\Psic+D\Zel\hat\eta+\rho$, where $\rho$ collects the coarse correction and the nonlinear part of the detail (\eqref{eq:eps2} and higher orders). Then
\begin{align}
\E\|x_1-x_0\|^2_{\rm phys}&=\E\|\rho\|^2,\\
\E\|x_1-x_0\|^2_{\rm indep}&=\E\|\rho+D\Zel(\hat\eta-\hat\eta')\|^2=\E\|\rho\|^2+2\,\E\|D\Zel\hat\eta\|^2+2\,\E\langle\rho,D\Zel\hat\eta\rangle,
\end{align}
with $\eta'$ an independent copy. The independent coupling adds twice the linear octave variance (plus a cross term that vanishes at second order because $\rho$ is quadratic in $\eta$): the network would have to ``re-invent'' the linear theory it already has. Under the physical coupling, $x_1-x_0$ vanishes at first order in perturbation theory, is small in single-stream regions, and is large only where the map $\eta\mapsto\Psif$ is strongly nonlinear (multi-stream regions). In the toy model the baseline $x_0$ alone already reaches $r\approx0.92$ against the truth in the octave band.

A remark on the geometry of the source. Conditional on $C$, $x_0-P\Psic$ is curl-free with covariance $\propto D^2P_\lin(k)k_ik_j/k^4$ per mode: a rank-one, purely longitudinal three-component field, i.e.\ a distribution supported on a $7N_c^3$-dimensional linear subspace of the $3\times7N_c^3$-dimensional detail space. This is not a defect. Conditional on $C$ the target is the push-forward of the \emph{same} scalar octave through $\Nbody_f$ (Proposition~\ref{prop:push}), hence also supported on a $7N_c^3$-dimensional (curved) manifold; the transverse components of the true detail are functions of $\eta$, not extra randomness. A regular flow maps one such manifold onto the other. What the degeneracy does imply is that the learned velocity field is only constrained on the support of $p_t$; off-manifold errors of the network are not self-correcting, which is an argument for keeping an Eulerian (density) term in the loss and for the multi-step integration, and it is one of the places where the idealised conditioning of Proposition~\ref{prop:push} and the reduced conditioning actually used differ.

\subsection{Emulator mode, generative mode, and the one-step (regression) limit}

\emph{Emulator mode} feeds the true octave $\eta$ of a simulation into \eqref{eq:x0} and integrates the learned ODE. \emph{Generative mode} feeds a sampled $\eta$; by Theorem~\ref{thm:fm} the output has the correct conditional law. In both modes the ODE is deterministic: all randomness of a sample enters through $x_0$, and the integration only transports it. Nothing is ``added back'' along the way; the marginal velocity $u_t$ is itself a conditional mean.

Consider a single explicit Euler step from $t=0$:
\begin{equation}
\hat x_1=x_0+u_0(x_0)=x_0+\E[x_1-x_0\mid x_0,C]=\E[x_1\mid x_0,C].
\end{equation}
The one-step model is exactly the conditional-mean regressor \emph{with respect to the information $(x_0,C)$}. What that implies depends on how much $(x_0,C)$ determines.

\paragraph{Full physical conditioning.} With the physical coupling, $x_0-P\Psic=D\Zel\hat\eta$ is the linear octave, so $\eta=-D^{-1}\nabla\cdot(x_0-P\Psic)$ is recovered exactly, and if $C$ determines $\delta_c$ then $x_1=\Nbody_f[\delta_c+\eta]$ is a deterministic function of $(x_0,C)$: $\Var(x_1\mid x_0,C)=0$. An ideal one-step regressor is then exact, and so is the ideal flow. Chaos does not alter this: a fixed initial condition has a single outcome. (In a real code, round-off amplified by chaotic dynamics makes the map effectively non-deterministic at the particle level inside halos; this is a statement about the simulator, not about the mathematics of the conditional.)

\paragraph{Partial information.} Whenever the conditioning leaves residual uncertainty --- coarse field only (no octave), the independent coupling, a projected or patch-local view of $C$, a finite-capacity network --- the law of total variance for any Fourier mode,
\begin{equation}
\E|\hat x_1|^2=\E\big|\E[\hat x_1\mid\mathcal I]\big|^2+\E\big[\Var(\hat x_1\mid\mathcal I)\big],
\end{equation}
says the conditional-mean predictor given the information $\mathcal I$ is deficient in power by the conditional variance. Because $\E\big[\hat x_1\,\overline{\E[\hat x_1\mid\mathcal I]}\big]=\E|\E[\hat x_1\mid\mathcal I]|^2$, its cross power with the truth equals its own power, hence
\begin{equation}
\boxed{\ \text{conditional mean given }\mathcal I:\quad \frac{P_{\rm pred}}{P_{\rm true}}=r^2 ;\qquad \text{correct generative model:}\quad \frac{P_{\rm pred}}{P_{\rm true}}=1\ \text{for any } r .\ }
\label{eq:r2}
\end{equation}
The relevant contrast is therefore between a regressor trained \emph{without} the octave (or with the independent coupling), which collapses to $\E[x_1\mid C]$ and sits on the line $P/P_{\rm true}=r^2$, and a generative model whose randomness is the octave, which sits on the line $P/P_{\rm true}=1$. With the physical coupling both objectives are in principle exact; the flow's advantage there is empirical (an iterative, easier regression at every $t$), not a variance argument.

\paragraph{Reading the toy numbers.} In the reference $16\to32$ run, one Euler step gives $r=0.993$ and $P_{\rm pred}/P_{\rm true}=0.959$, whereas an optimal conditional mean with that $r$ would have $r^2=0.986$. The one-step model is therefore \emph{not} the optimal regressor; the gap is a model/optimisation deficit, and the improvement to $0.993$ with eight Heun steps is best described as the multi-step parameterisation fitting the same map better, to be checked against a direct regression baseline trained under identical conditions. The sampled-octave run has $P_{\rm pred}/P_{\rm true}\approx1.02$ with $r\approx0$, as \eqref{eq:r2} predicts for a generative model.

\begin{exercise}
Prove the identity $\E[\hat x_1\overline{\E[\hat x_1\mid\mathcal I]}]=\E|\E[\hat x_1\mid\mathcal I]|^2$ and deduce $r^2=P_{\rm pred}/P_{\rm true}$ for the conditional-mean predictor. Then show that for the physical coupling with $C$ determining $\delta_c$, $\Var(x_1\mid x_0,C)=0$.
\end{exercise}

\subsection{Sampling}

The ODE $\dot X_t=v_\theta(X_t,t)$ is integrated from $X_0=x_0$ with Heun's method on a uniform grid of $n$ steps. Because $v_\theta$ is trained on straight-line interpolants under a short coupling, few steps are needed; the sweep $n=1,2,8$ gives octave-band power ratios $0.959,0.978,0.993$ on the reference model, which by the discussion above measures how well the multi-step parameterisation fits the map, not a restoration of variance. The wavelet-conditional argument of Guth et al.\ (2022) is that, when generation is organised scale by scale, the number of steps per scale is independent of resolution; here each $2\times$ step always generates ``the next octave relative to its input'', so the task is the same at every level.

\subsection{Sampling the linear octave}

To draw $\eta$ with the right statistics without an external linear power spectrum, the code measures the isotropic power of the linear \emph{density} from the training initial conditions, $\hat\delta^{(1)}=-i\kk\cdot\hat\Psi^{(1)}$, and synthesises a Gaussian $\hat\delta$ with that power on the modes outside $W$ (with the cube window these are exactly the octave; with the sphere they also include the corner modes, subject to the caveat of Remark~\ref{rem:clean}); the displacement then follows from \eqref{eq:zel}, which makes it curl-free automatically. Sampling the three displacement components independently would not be curl-free and would be wrong at first order.

%======================================================================
\section{Scale conditioning and self-similarity}\label{sec:ss}
%======================================================================

\subsection{Dimensionless inputs}

The displacement has units of length and its large-scale part (bulk flows) is unbounded in units of $h_\ell$ as $\ell$ grows. The deformation tensor $D_{ij}=\partial_i\Psi_j$ is dimensionless, invariant under translations of $\Psi$, and its statistics on resolved scales do not depend on the grid. The detail in grid units, $d/h_f$, has a bounded amplitude controlled by the linear variance of the octave: with $P_\Psi\propto P_\lin/k^2$, the octave variance of the displacement scales as $\sigma_{\Psi,\ell}^2\propto k_\ell^{\,n+1}$ and $\sigma_{\Psi,\ell}/h_\ell\propto k_\ell^{(n+3)/2}$, a slowly varying quantity for $n>-3$. A single network acting on $(x_t-P\Psic)/h_f$ and $D_{ij}$ therefore sees inputs of the same nature at every level, and only a scalar ``how nonlinear is this level'' needs to be supplied.

\subsection{Exact self-similarity in a scale-free cosmology}

For $P_\lin(k)=Ak^n$ in an Einstein--de Sitter background the continuum dynamics has no intrinsic scale. The linear variance smoothed on scale $R$ is
\begin{equation}
\sigma^2(R,a)=D^2(a)\int\frac{\dd^3k}{(2\pi)^3}Ak^n\,\big|\tilde W(kR)\big|^2\propto a^2R^{-(n+3)} ,
\end{equation}
so the nonlinear scale defined by $\sigma(R_{\rm NL},a)=1$ is $R_{\rm NL}(a)\propto a^{2/(n+3)}$, and every dimensionless statistic of the evolved field depends on a length $r$ only through $r/R_{\rm NL}(a)$. For the discretised problem with grid spacing $h_\ell$ (and force softening $\propto h_\ell$) the dimensionless state of level $\ell$ at time $a$ is $h_\ell/R_{\rm NL}(a)$.

\begin{proposition}[Level shift equals time shift]
The $2\times$ step from level $\ell+1$ at expansion factor $a$ is statistically identical, for local statistics, to the step from level $\ell$ at $a'$ with $h_\ell/R_{\rm NL}(a')=h_{\ell+1}/R_{\rm NL}(a)$, i.e.
\begin{equation}
R_{\rm NL}(a')=2R_{\rm NL}(a)\quad\Longleftrightarrow\quad a'=2^{(n+3)/2}\,a .
\end{equation}
For $n=-2$, $a'=\sqrt2\,a$.
\end{proposition}
Going one level finer at fixed time is the same as waiting until $R_{\rm NL}$ has doubled at fixed grid. Finite-box effects break this only for the longest modes, which is why the statement is about local statistics; the operator is local, so that is enough. This gives a strict test of weight sharing and free data augmentation in scale-free simulations.

\subsection{The universality hypothesis in $\Lambda$CDM}

Without exact self-similarity, condition the shared operator on
\begin{equation}
s_\ell=\big(\sigma(h_\ell,z),\ n_{\rm eff}(k_\ell),\ f(z)\big),
\end{equation}
the linear rms fluctuation at the grid scale, the local slope of the linear spectrum there, and the growth rate (which fixes the velocity--displacement relation). The hypothesis, analogous to the universality of the halo mass function in $\nu=\delta_c/\sigma$, is that the $2\times$ operator depends on $(\ell,z,\text{cosmology})$ only through $s_\ell$. It is testable on the $64/128/256/512$ series by comparing transitions with matched $\sigma$.

%======================================================================
\section{Evaluation metrics}
%======================================================================

All spectra are shell averages of $|\hat a|^2$, $|\hat b|^2$ and $\mathrm{Re}(\hat a\hat b^*)$ summed over the three components, normalised as in Section~1.2, and computed separately in the coarse band ($W$) and the octave band ($1-W$).

\begin{itemize}
\item[] \emph{Cross-correlation coefficient} $r(k)=P_{ab}/\sqrt{P_{aa}P_{bb}}$ between the model output $a$ and the truth $b$. In emulator mode it measures determinism; $1-r^2$ is the fraction of power a generative model must supply.
\item[] \emph{Power ratio} $P_{aa}/P_{bb}$: must be $1$ in generative mode; the pair $(r,P_{aa}/P_{bb})$ locates a model between the regression line $P_{aa}/P_{bb}=r^2$ and the generative line $P_{aa}/P_{bb}=1$, cf.\ \eqref{eq:r2}.
\item[] \emph{Relative correction error} in the coarse band, $\E|\hat a-\hat b|^2/\E|\hat b|^2=(P_{aa}+P_{bb}-2P_{ab})/P_{bb}$: how well the step repairs the coarse run. Its generative-mode value exceeds its emulator-mode value because the $\eta\eta$ backreaction of Proposition~\ref{prop:k2} depends on the octave realisation.
\item[] \emph{Wiener transfer function} $T(k)=P_{cf}/P_{ff}$ between the prolonged coarse run and the fine run. It is the isotropic linear filter that minimises $\E|T\Psif-\Psic|^2$ (standard Wiener argument: minimise a quadratic in $T$ mode by mode), i.e.\ the coarse-graining a real low-resolution run most resembles; it decides where $\alpha$ should be placed.
\end{itemize}

%======================================================================
\section{What is proven, what is assumed, what is measured}
%======================================================================

Proven above: the operator identities of Section~2; Proposition~\ref{prop:linear}; the split \eqref{eq:eps2} including the coarse-only harmonics in the detail, Lemma~\ref{lem:kernel} and the ensemble form of Proposition~\ref{prop:k2} for the octave--octave term; the locality statement \eqref{eq:B}; the push-forward statement in its mixture form \eqref{eq:mixture} and the Markov property of the full nested chain; Theorem~\ref{thm:fm} and Proposition~\ref{prop:loss}; the information-set form of the regression identity \eqref{eq:r2}; the scale-free level--time correspondence.

Idealisations to keep in view: the push-forward collapses to a single map only when the conditioning determines $\delta_c$; the reduced states actually used (projected fields, patches, finite networks) make it a mixture, and make the reduced chain Markov only approximately; with the sphere window the corner modes are re-sampled although the retained field is informative about them (the cube window avoids this); the $k^2$/$k^4$ law is for the ideal second-order octave--octave term, the mixed term is at least as suppressed, and neither higher orders nor discreteness errors are claimed to share the exponent; the ``$\vartheta$'' spectra are divergences of displacements, not Eulerian densities.

Changes relative to the first version of these notes, following an external review: the real-space sign of $\Psi^{(2)}$ in \eqref{eq:psi2}; the missing $S[\delta_L]$ term in $d^{(2)}$ of \eqref{eq:eps2}; divergence versus density in Proposition~\ref{prop:k2} and Remark~\ref{rem:density}; the Wick-theorem form of the $k^4$ statement; the mixture form of the conditional and the scope of the Markov property; the removal of the claim that one-step regression necessarily loses power under full physical conditioning and of the claim that the multi-step flow ``adds back'' variance; and the recommendation of the cube window as the theoretical default.

Assumed, to be tested on the real resolution series: that the discreteness floor of a real low-resolution run is small (panel (a) of Phase~0); that the conditional of the detail given the local deformation tensor is close to Gaussian outside multi-stream regions and still tractable inside them (panel (c)); that the universality in $s_\ell$ holds well enough for weight sharing (matched-$\sigma$ comparison). The toy experiments use 2LPT, which has no shell crossing; nothing in them speaks to the multi-stream regime.

\section*{References}
\begin{itemize}
\item[] Bouchet, Colombi, Hivon, Juszkiewicz 1995, A\&A 296, 575 (2LPT conventions).
\item[] Li, Ni, Croft, Di Matteo, Bird, Feng 2021, PNAS (AI-assisted SR I); Zhang et al.\ 2024, MNRAS 528, 281 (III); Zhang et al.\ 2024, OJAp (IV).
\item[] Lipman, Chen, Ben-Hamu, Nickel, Le 2023, ICLR (flow matching); Tong et al.\ 2023 and Pooladian et al.\ 2023 (arbitrary couplings / minibatch OT); Albergo \& Vanden-Eijnden 2023 (stochastic interpolants).
\item[] Guth, Coste, De Bortoli, Mallat 2022, NeurIPS, arXiv:2208.05003 (wavelet score-based generative modelling); Marchand, Ozawa, Biroli, Mallat 2022, arXiv:2207.04941; Guth, Lempereur, Bruna, Mallat 2023, arXiv:2306.00181.
\item[] Schmidt, Cabass, Jasche, Lavaux 2018, arXiv:1803.03274; St\"ucker et al.\ 2021, MNRAS 503, 1473 (anisotropic separate universe).
\item[] Joyce, Garrison, Eisenstein 2021, MNRAS 504, 3550 (scale-free self-similarity).
\item[] Nishimichi, Bernardeau, Taruya 2016, PLB, arXiv:1411.2970 (response of large scales to small scales).
\end{itemize}

\end{document}
